\section{Rational Inequalities and Applications}

\subsection{Variation}

\begin{definition}[Variations]
    Suppose $x$, $y$ and $z$ are variable quantities. We say
    \begin{itemize}
        \item $y$ \textbf{varies directly with} (or is \textbf{directly proportional to}) $x$ if there is a constant $k$ such that $y=kx$
        \item $y$ \textbf{varies inversely with} (or is \textbf{inversely proportional to}) $x$ if there is a constant $k$ such that $y=\frac{k}{x}$.
        \item $z$ \textbf{varies jointly with} (or is \textbf{jointly proportional to}) $x$ and $y$ if there is a constant $k$ such that $z=kxy$.
    \end{itemize}
    The constant $k$ in the above definitions is called the \textbf{constant of proportionality}
\end{definition}

\subsection{Exercises}
\includepdf[pages=-]{Exercises/04/03.pdf}